\documentclass[12pt,a4paper]{article}
\usepackage[utf8]{inputenc}
\usepackage[portuguese]{babel}
\usepackage{amsmath}
\usepackage{amsfonts}
\usepackage{amssymb}
\usepackage{mathrsfs}
\usepackage{graphicx}
\usepackage{float}
\usepackage[left=3cm,right=2cm,top=3cm,bottom=2cm]{geometry}
\usepackage{pgfplots}
\pgfplotsset{compat=1.17}
\usepackage{tikz}
\usepackage{booktabs}
\usepackage{siunitx}
\usepackage{caption}
\usepackage{subcaption}
\usepackage{indentfirst}
\usepackage{hyperref}
\usepackage{xurl}

\hypersetup{
    colorlinks=true,
    linkcolor=black,
    citecolor=black,
    urlcolor=blue
}

\renewcommand{\baselinestretch}{1.5}

\title{
    \vspace{2cm}
    \Huge\textbf{\MakeUppercase{Predição de Mutagenicidade Ames}}\\[0.5cm]
    \Large Trabalho Prático 1 - INF01017\\[0.3cm]
    \large Aprendizado de Máquina\\[0.3cm]
    \large Instituto de Informática --- UFRGS\\[0.5cm]
    \vspace{2cm}
}

\author{
    \textbf{Professora:} Mariana Recamonde Mendoza\\[2em]
    \textbf{Alunos:}\\
    Everton Fritsch de Lima --- 00334081\\
    [Nome Integrante 2] --- [Cartão]\\
    [Nome Integrante 3] --- [Cartão]\\
}

\date{
    \vfill
    Porto Alegre\\
    Maio de 2026
}

\begin{document}
\maketitle

\section{Definição do Problema e Coleta de Dados}

% Descreva aqui:
% - O problema de classificação/regressão escolhido
% - Dataset utilizado (link, referência)
% - Objetivo do trabalho
% - Descrição dos dados (número de instâncias, atributos, target)

\section{Análise Exploratória e Pré-processamento dos Dados}

% Descreva aqui:
% - Características do dataset (tipos de atributos, distribuições)
% - Análise exploratória (gráficos, estatísticas)
% - Problemas identificados (valores faltantes, desbalanceamento, etc.)
% - Técnicas de pré-processamento aplicadas
% - Justificativa das escolhas de pré-processamento

\section{Abordagem, Algoritmos e Estratégia de Avaliação}

% Descreva aqui:
% - Abordagem escolhida (classificação ou regressão)
% - Algoritmos selecionados (mínimo 5, diversificados)
% - Justificativa da diversidade dos algoritmos
% - Métricas de avaliação escolhidas
% - Estratégia de validação (holdout, cross-validation, etc.)
% - Configurações de hiperparâmetros (se aplicável)

\section{Spot-checking de Algoritmos}

% Descreva aqui:
% - Metodologia de execução dos experimentos
% - Resultados obtidos (tabelas, gráficos)
% - Análise comparativa do desempenho dos modelos
% - Identificação dos 2-3 algoritmos mais promissores
% - Discussão dos resultados

\section{Conclusão}

% Sumarize:
% - Principais achados do trabalho
% - Algoritmos mais promissores identificados
% - Próximos passos (fora do escopo do T1)

\end{document}